Расстояние Громова--Хаусдорфа, рассматриваемое в данной работе,
является расстоянием
между метрическими пространствами. Впервые оно было введено в \cite{Edwards}.

Традиционно, свойства расстояния
Громова-Хаусдорфа
изучаются на множестве компактных метрических пространств,
рассматриваемых с точностью до изометрии. На данном множестве
расстояние удовлетворяет определению метрики. Именно такое расстояние
рассматривал М. Громов в \cite{gromov_structures_1981}, где получил
некоторые фундаментальные свойства данной метрики. Так, одним из
важнейших результатов стал тот факт, что метрическое пространство
компактов является полным и стягиваемым. Под стягиваемостью здесь
подразумевается наличие отображения гомотетии, то есть умножения
пространств на положительное число \( \lambda  \). Данное умножение
является непрерывным с точки зрения метрики Громова-Хаусдорфа и при
стремлении \( \lambda  \) к \( 0 \) все компактные метрические
пространства \( \lambda X \) сходятся к одноточечному пространству \(
\Delta _1 \).

Позднее, в \cite{gromov_metric_2001} М. Громов рассматривал
расстояние Громова-Хаусдорфа на классе всех, не обязательно
ограниченных метрических пространств. В этом случае расстояние не
удовлетворяет всем аксиомам метрики, в частности расстояние между
двумя неограниченными метрическими пространствами может быть равно
бесконечности. В этом случае оказалось удобным провести факторизацию
всего пространства по конечным расстояниям, и рассматривать отдельных
представителей полученного факторпространства. Такие классы в
дальнейшем назвали \emph{облаками}. То есть облако --- все
метрические пространства, находящиеся на конечном расстоянии
Громова--Хаусдорфа от данного. Здесь и далее, облако содержащее
метрическое пространство \( X \) будем обозначать \( [X] \).

Облака обладают многими нетривиальными свойствами. Так, прежде, чем
говорить об их геометрии, стоит рассмотреть облака с точки зрения
теории множеств. Традиционно, в работах, связанных с ними
используется система аксиом NBG (\cite{bernays_system_1937},
\cite{godel_consistency_1940}, \cite{neumann_1925}). В ней
объектами изучения являются \emph{классы}. Классы, которые являются
элементами других классов полагаются множествами. Остальные классы
называются \emph{собственными}. В данной работе показано, что облака
не являются множествами, то есть являются собственными классами.
Собственный класс, снабженный метрикой называется метрическим
классом. Чтобы каждое облако можно было считать метрическим классом,
в данной работе отождествляются метрические пространства, находящиеся
на нулевом расстоянии Громова--Хаусдорфа друг от друга. Вместе с
этим, можно рассматривать расстояние Громова--Хаусдорфа между
облаками, полностью аналогичное расстоянию между метрическими
пространствами. Исследованию его свойств посвящена основная часть данной работы.

В работе \cite{gromov_metric_2001} М. Громов предполагал, что все
облака, аналогично пространству ограниченных метрических пространств,
являются полными и стягиваемыми. Полнота облаков была доказана
в работе \cite{TuzhBog1}.

Для изучения проблемы стягиваемости на каждом облаке вводится
операция умножения на положительное число. Более точно, образ облака
\([X] \) при умножении на \( \lambda \in (0, \infty) \) --- облако \(
\lambda [X] \), состоящее из всех пространств из облака \( [X] \),
умноженных на \( \lambda \). Если предполагать, что все облака
стягиваемы, то должны выполняться следующие условия:
\begin{enumerate}
  \item облака \( [X] \) и \( \lambda [X] \) должны
    быть равны,
  \item все пространства \( \lambda Y \) должны непрерывно стремиться
    к одному и тому же пространству в облаке \( [X] \).
\end{enumerate}
Оказалось, что в общем случае оба этих утверждения не верны.

Можно привести пример неограниченного метрического пространства \( X
\), для которого расстояние между \( X \) и \( \lambda X \) равно
бесконечности для некоторых \( \lambda \in \R^+ \)(см. например
\cite{TuzhBog1}). Тогда все все пространства из облака \( [X] \ni X\)
перейдут в облако \( [\lambda X] \ni \lambda X \), не совпадающее с
\( [X] \). В работе \cite{TuzhBog2} было показано, что те \( \lambda
\), для которых облака \( [X] \) и \( [\lambda X] \) совпадают
образуют мультипликативную подгруппу в \( \R^+ \) для фиксированного
облака \( [X] \). Такая подгруппа называется \emph{стационарной
группой} облака \( [X] \).

В данной работе я исследую свойства облаков, обладающих
нетривиальными стационарными группами. Основное внимание уделено
парам облаков, стационарные группы которых имеют пересечение,
отличное от \( \{1\} \).
